\begin{definition}[The gather-based GBCA implementation]
  \label{def:aba-round-over-bracha}
  \lean{PLTS.ABA.GBCA.ByAFW.RoundStateOverBroadcastSpecification}\lean{PLTS.ABA.GBCA.ByAFW.roundOverBroadcastSpecification}\lean{PLTS.ABA.GBCA.ByAFW.RoundStateOverBracha}\lean{PLTS.ABA.GBCA.ByAFW.roundOverBracha}\leanok
  \uses{def:aba-round-composition, def:aba-round-over-gather-specifications, def:aba-gather-over-broadcast-specification, def:aba-gather-over-bracha}
  The round of Definition~\ref{def:aba-round-composition} at two concrete gather systems in its two gather positions.
  The composition over the broadcast specifications places two instances of Definition~\ref{def:aba-gather-over-broadcast-specification} in those positions; the fully concrete composition places two instances of Definition~\ref{def:aba-gather-over-bracha} --- the two gather instances and, beneath them, the $4n$ Bracha instances carrying the inputs and the $\mathrm{BIND}$ payloads.
  In both, the round's $\mathrm{callG}$ is the first gather's call, taken by the caller's program and by the first gather at once; the first gather's return to a process and that process's call of the second gather at the candidate are two hidden events, the second recording the candidate as the caller's gather input, and the call of the caller's input-broadcast instance is a separate internal transition of that gather; the second gather's return and the round's $\mathrm{retG}$ are two events likewise; each component's internal transitions are silent transitions of the round; and corruption is the family's broadcast into the two gathers (D1).
  The concrete composition \leandecl{PLTS.ABA.GBCA.ByAFW.roundOverBracha} is the gather-based GBCA implementation, the protocol as it runs; the composition over the broadcast specifications, \leandecl{PLTS.ABA.GBCA.ByAFW.roundOverBroadcastSpecification}, is the one the two substitutions below pass through.
\end{definition}
